% Standalone problem document, generated from the adjacent README.md.
% Compile with XeLaTeX. Mathematical content is maintained in Markdown.
\documentclass[11pt,a4paper]{article}
\usepackage[fontsize=10.5pt]{scrextend}
\usepackage[margin=22mm,headheight=15pt,headsep=9mm,footskip=11mm]{geometry}
\usepackage{fontspec}
\setmainfont{texgyrepagella}[Extension=.otf,UprightFont=*-regular,BoldFont=*-bold,ItalicFont=*-italic,BoldItalicFont=*-bolditalic]
\setsansfont{texgyreheros}[Extension=.otf,UprightFont=*-regular,BoldFont=*-bold,ItalicFont=*-italic,BoldItalicFont=*-bolditalic]
\setmonofont{texgyrecursor}[Extension=.otf,UprightFont=*-regular,BoldFont=*-bold,ItalicFont=*-italic,BoldItalicFont=*-bolditalic,Scale=MatchLowercase]
\usepackage{amsmath,amssymb}
\usepackage{unicode-math}
\setmathfont{texgyrepagella-math.otf}
\usepackage{xcolor,microtype,enumitem,needspace,fancyhdr}
\usepackage{bookmark}
\definecolor{ink}{HTML}{17344B}
\definecolor{accent}{HTML}{197D80}
\definecolor{muted}{HTML}{596773}
\hypersetup{colorlinks=true,linkcolor=accent,urlcolor=accent,pdftitle={RA-05 - Sharp
joint rank and accuracy dependence for strong \ell\_p subspace
coresets},pdfauthor={Open Problems in Numerical Linear Algebra}}
\urlstyle{same}
\setlength{\parindent}{0pt}
\setlength{\parskip}{5pt plus 1pt minus 1pt}
\linespread{1.04}
\setlength{\emergencystretch}{3em}
\widowpenalty=10000
\clubpenalty=10000
\displaywidowpenalty=10000
\allowdisplaybreaks[1]
\setlist{nosep,leftmargin=1.5em}
\providecommand{\tightlist}{\setlength{\itemsep}{0pt}\setlength{\parskip}{0pt}}
\providecommand{\pandocbounded}[1]{#1}
\pagestyle{fancy}
\fancyhf{}
\fancyhead[L]{\sffamily\footnotesize\color{muted}OPEN PROBLEMS IN NUMERICAL LINEAR ALGEBRA}
\fancyhead[R]{\sffamily\bfseries\footnotesize\color{accent}RA-05}
\fancyfoot[L]{\parbox[b]{0.9\textwidth}{\sffamily\fontsize{8}{10}\selectfont\color{muted}Editorial ratings; a bounded search cannot certify that a problem remains unsolved.\\Literature check: 2026-09-14}}
\fancyfoot[R]{\sffamily\footnotesize\color{muted}\thepage}
\renewcommand{\headrulewidth}{0.3pt}
\renewcommand{\headrule}{\hbox to\headwidth{\color{accent}\leaders\hrule height \headrulewidth\hfill}}
\makeatletter
\renewcommand{\section}{\Needspace{5\baselineskip}\@startsection{section}{1}{0pt}{12pt plus 2pt minus 2pt}{4pt}{\sffamily\bfseries\large\color{ink}}}
\renewcommand{\subsection}{\Needspace{4\baselineskip}\@startsection{subsection}{2}{0pt}{9pt plus 1pt minus 1pt}{3pt}{\sffamily\bfseries\normalsize\color{ink}}}
\makeatother
\setcounter{secnumdepth}{0}
\begin{document}
{\sffamily\small\color{muted}Randomized and low-rank approximation\par}
\vspace{3pt}
{\raggedright\hyphenpenalty=10000\exhyphenpenalty=10000\sffamily\bfseries\fontsize{21}{24}\selectfont\color{ink}Sharp
joint rank and accuracy dependence for strong \(\ell_p\) subspace
coresets\par}
\vspace{7pt}
{\sffamily\small\textbf{Difficulty:} challenging\quad\textbf{Importance:} interesting
to the community\par}
{\sffamily\small\bfseries\color{accent}Status: Solved\par}
\vspace{4pt}
{\color{accent}\hrule height 0.8pt}
\vspace{6pt}
\textbf{Rating rationale:} Challenging because simultaneous control of
all subspaces must match joint rank and accuracy lower bounds; community
impact is compact robust low-rank fitting.

\subsection{Full resolution --- even/non-even classification, 14
September
2026}\label{full-resolution-evennon-even-classification-14-september-2026}

\textbf{Author:} Sidney Holden, Center for Computational Biology,
Flatiron Institute, Simons Foundation.
\href{https://github.com/ajt60gaibb/OpenProblemsInNLA/blob/main/references/holden-further-2026-09-14/README.md}{Verified
affiliation and submission record}.

\href{https://github.com/ajt60gaibb/OpenProblemsInNLA/blob/main/references/holden-further-2026-09-14/RA-05/manuscript.pdf}{Part
I, Theorem 1.1, and Part II, Theorem 1.1} classify the optimal
worst-case strong original-row coreset size for every fixed real
\(p>2\), all \(k\ge1\) and \(0<\varepsilon<1/2\), with arbitrary input
rank and ambient dimension. Writing \(S_p(k,\varepsilon)\) for that
size, the result is

\[S_p(k,\varepsilon)=\widetilde\Theta_p\!\left(
\begin{cases}
k^{p/2}\varepsilon^{-2},&p\notin2\mathbb Z,\\
\min\{k^{p/2}\varepsilon^{-2},\ k^{(p+1)/2}\varepsilon^{-1}+k^{p/2-1}\varepsilon^{-2}\},&p\in2\mathbb Z.
\end{cases}\right)\]

The hidden factors are logarithmic in \(2k/\varepsilon\); constants may
depend only on the fixed exponent. Weights are nonnegative and supported
on original rows, and the guarantee holds simultaneously for every
fitted subspace of dimension at most \(k\). The displayed additive
proposal in the retained problem is false for every fixed \(p>2\). This
completes the broader classification as well as its subsidiary truth
question. The general upper bound for non-even powers is credited to
Lin--Mirrokni--Woodruff; the manuscript credits Li--Wang--Woodruff for
the antecedent Fourier mechanism.

Both complete proof parts passed a separate
\href{https://github.com/ajt60gaibb/OpenProblemsInNLA/blob/main/reviews/2026-09-14-further-submissions/RA-05-review.md}{independent
Codex AI-agent informal audit}, including the new uniform estimates and
primary-source dependency checks. This meets the repository's
\textbf{Solved} policy. The exact cubic certificate passed 17,575
assertions; optional numerical diagnostics were not rerun because their
dependencies were unavailable. The proof audit, rather than those finite
checks, supports the universal result.
\href{https://github.com/ajt60gaibb/OpenProblemsInNLA/blob/main/references/holden-further-2026-09-14/RA-05/part-1.tex}{Part
I source} ·
\href{https://github.com/ajt60gaibb/OpenProblemsInNLA/blob/main/references/holden-further-2026-09-14/RA-05/part-2.tex}{Part
II source}.

ChatGPT assistance is disclosed. No Lean verification, external human
peer review, formal verification or priority certification is asserted.
The earlier partial results and historical open-status checks below are
retained as history; this dated full resolution supersedes their
remaining-gap statements. Original ID, canonical path and mathematical
target are unchanged.

\subsection{Further partial resolution --- even-power high accuracy, 13
September
2026}\label{further-partial-resolution-even-power-high-accuracy-13-september-2026}

\textbf{Author:} Sidney Holden, Center for Computational Biology,
Flatiron Institute, Simons Foundation.
\href{https://github.com/ajt60gaibb/OpenProblemsInNLA/blob/main/references/holden-ra05-even-power-2026-09-13/README.md}{Submission
and verified affiliation}.

For each fixed integer \(s\ge2\),
\href{https://github.com/ajt60gaibb/OpenProblemsInNLA/blob/main/references/holden-ra05-even-power-2026-09-13/manuscript.pdf}{Theorems
1.1--1.2 and Corollary 1.3} establish

\[c_s\frac{k^{s-1}}{\varepsilon^2}\le S_{2s}(k,\varepsilon)
\le C_s\left(\frac{k^{(3s-1)/2}}{\varepsilon}+\frac{k^{s-1}}{\varepsilon^2}\right)
\log^5\!\left(\frac{2k}{\varepsilon}\right).\]

Thus \(S_{2s}(k,\varepsilon)=\widetilde\Theta_s(k^{s-1}/\varepsilon^2)\)
whenever \(0<\varepsilon<1/2\) and \(\varepsilon\le k^{-(s+1)/2}\). The
bounds hold for every \(k\ge1\), with arbitrary input rank and ambient
dimension, nonnegative original-row weights, and simultaneous
preservation of every subspace of dimension at most \(k\). Constants
depend only on \(s\); no efficient construction is claimed.

The stated partial results passed an
\href{https://github.com/ajt60gaibb/OpenProblemsInNLA/blob/main/references/holden-ra05-even-power-2026-09-13/verification/independent-review.md}{independent
Codex AI-agent informal audit}, including primary-source dependency
checks and 8,453 finite assertions. Finite checks are supporting
diagnostics, not proofs of the asymptotic theorem. ChatGPT assistance is
disclosed; no external human peer review or formal verification is
asserted. No Lean verification was performed.

\textbf{RA-05 remains Partially resolved.} Non-even exponents and
intermediate accuracies for even \(p\ge6\) remain unclassified. The
earlier quartic classification and all-exponent lower bounds below are
retained as separate prior contributions.
\href{https://github.com/ajt60gaibb/OpenProblemsInNLA/blob/main/references/holden-ra05-even-power-2026-09-13/manuscript.tex}{New
proof source}.

\subsection{Partial resolution --- unrestricted quartic case, 13
September
2026}\label{partial-resolution-unrestricted-quartic-case-13-september-2026}

\textbf{Author:} Sidney Holden, Center for Computational Biology,
Flatiron Institute, Simons Foundation.
\href{https://github.com/ajt60gaibb/OpenProblemsInNLA/blob/main/references/holden-ra05-quartic-2026-09-13/README.md}{Submission
and verified affiliation}.

\href{https://github.com/ajt60gaibb/OpenProblemsInNLA/blob/main/references/holden-ra05-quartic-2026-09-13/manuscript/RA05_unrestricted_quartic.pdf}{Theorem
1.1, proved in Sections 2--8} classifies the optimal size for \(p=4\),
arbitrary input rank, every integer \(k\ge1\) and every
\(0<\varepsilon<1/2\), up to logarithmic factors:

\[S_4(k,\varepsilon)=\widetilde\Theta\!\left(
\min\left\{\frac{k^2}{\varepsilon^2},
\frac{k^{5/2}}{\varepsilon}+\frac{k}{\varepsilon^2}\right\}\right).\]

The weights are nonnegative and select original rows; preservation is
simultaneous over all subspaces of dimension at most \(k\), without a
restriction on input rank, row count or ambient dimension. The lower
bound has no logarithmic loss, the combined upper bound uses at most
\(\log^9(2k/\varepsilon)\), and the second upper branch uses at most the
fifth power.
\href{https://github.com/ajt60gaibb/OpenProblemsInNLA/blob/main/references/holden-ra05-quartic-2026-09-13/manuscript/RA05_unrestricted_quartic.tex}{Proof
source}.

This also disproves the subsidiary displayed upper-bound conjecture
below at \(p=4\): at \(\varepsilon=k^{-1}\) the lower bound is of order
\(k^{7/2}\), exceeding its proposed order \(k^3\) times any fixed
logarithmic power. \textbf{The unrestricted optimal joint classification
for fixed real \(p>2\) other than \(4\) remains open.} The whole RA-05
entry is therefore Partially resolved, and remains in the open count.

The argument passed a separate
\href{https://github.com/ajt60gaibb/OpenProblemsInNLA/blob/main/references/holden-ra05-quartic-2026-09-13/verification/independent-review.md}{independent
Codex AI-agent informal audit}, including primary-source checks of the
imported theorems. All 7,621 finite assertions and both exact
certificate checkers were rerun successfully. ChatGPT assistance is
disclosed; this is not external human peer review or formal
verification. No Lean verification was performed. The
\href{https://github.com/ajt60gaibb/OpenProblemsInNLA/pull/199}{earlier
all-exponent lower-bound submission, PR \#199}, is a separate partial
contribution. The original statement, ID, path and prior-source credit
are retained below.

\subsection{All-exponent lower bound --- 13 September
2026}\label{all-exponent-lower-bound-13-september-2026}

Holden's separate
\href{https://github.com/ajt60gaibb/OpenProblemsInNLA/blob/main/references/holden-ra05-2026-09-13/manuscript/RA05_all_p_lower_bounds.pdf}{Theorem
1.1} proves, for every fixed real \(p>2\), all sufficiently large \(k\)
and every \(0<\varepsilon<1/2\),

\[S_p(k,\varepsilon)\ge c_p\frac{k^{p/2}}{\varepsilon^{\beta_p}+(\log k)/k},
\qquad
\beta_p=\begin{cases}2,&p\in\{4,6,8,\ldots\},\\2-2/p,&\text{otherwise}.\end{cases}\]

One input per rank, independent of accuracy, witnesses the bound for
arbitrary nonnegative original-row weights. Corollary 1.2 disproves the
subsidiary additive formula for \textbf{every fixed real exponent
greater than two}, even allowing any fixed logarithmic power. Corollary
1.3 matches the cited upper bound up to logarithms for even exponents at
least four when \(\varepsilon\ge\sqrt{(\log k)/k}\). The
\href{https://github.com/ajt60gaibb/OpenProblemsInNLA/blob/main/references/holden-ra05-2026-09-13/verification/independent-review.md}{independent
review} and
\href{https://github.com/ajt60gaibb/OpenProblemsInNLA/blob/main/references/holden-ra05-2026-09-13/README.md}{submission
record} are retained. Together with the quartic classification above,
these results leave the other exponents' unrestricted optimal joint
sizes unresolved.

\subsection{Original question}\label{original-question}

Fix a real \(p>2\). For \(A\in\mathbb R^{n\times d}\), with rows
\(a_i^T\), and a linear subspace \(F\subseteq\mathbb R^d\), define

\[\mathop{\mathrm{cost}}\nolimits_A(F)=\sum_{i=1}^n\|a_i^T(I-P_F)\|_2^p,\]

where \(P_F\) is the Euclidean orthogonal projector. A strong row
coreset consists of nonnegative weights \(w_1,\ldots,w_n\), supported on
a small subset of the original rows, for which

\[(1-\varepsilon)\mathop{\mathrm{cost}}\nolimits_A(F)
\leq\sum_iw_i\|a_i^T(I-P_F)\|_2^p
\leq(1+\varepsilon)\mathop{\mathrm{cost}}\nolimits_A(F)\]

holds simultaneously for every \(F\) of dimension at most \(k\).

Determine the optimal worst-case coreset size as a joint function of
\(k\) and \(\varepsilon\), up to logarithmic factors. In particular, do
constants \(C_p,c_p>0\) exist such that, for all \(n,d\), integers
\(1\leq k< d\), all input matrices \(A\), and all \(0<\varepsilon<1/2\),
such a coreset exists with

\[|\mathop{\mathrm{supp}}\nolimits(w)|\leq
C_p\left(\frac{k^{p/2}}{\varepsilon}+\frac{k}{\varepsilon^2}\right)
\log^{c_p}\!\left(\frac{2k}{\varepsilon}\right)?\]

The constants may depend on fixed \(p\), but not on
\(n,d,k,\varepsilon\). This is an existence/size question; no extra
running-time target is imposed.

Such a coreset permits subsequent robust low-rank fitting on a much
smaller weighted matrix. The August 2026 paper proves an upper bound of
order \(\widetilde O_p(k^{p/2}/\varepsilon^2)\) and explicitly
identifies the remaining gap to the joint lower-bound expression in
Section 1.4. Its new result supersedes older questions that merely asked
for an \(\varepsilon^{-2}\) dependence.

\subsection{References}\label{references}

\begin{enumerate}
\def\labelenumi{\arabic{enumi}.}
\tightlist
\item
  H. Lin, V. Mirrokni, and D. P. Woodruff, \emph{Nearly Optimal Strong
  Coresets for \(\ell_p\) Subspace Approximation}, arXiv:2608.26047v2,
  Section 1.4 ``Open problems''; Section 1.2 for the new bounds.
  \href{https://arxiv.org/html/2608.26047v2}{Paper}.
\item
  D. P. Woodruff and T. Yasuda, \emph{Root ridge leverage score sampling
  for \(\ell_p\) subspace approximation}, FOCS 2025, full version
  arXiv:2407.03262. \href{https://arxiv.org/abs/2407.03262}{Paper}.
\end{enumerate}

\subsection{Status check --- 2026-09-08}\label{status-check-2026-09-08}

Searched ``strong coresets subspace approximation p greater than 2
epsilon 2026'' and the exact 2026 title; checked the latest arXiv v2 of
August 27, 2026. No later resolution was found. The target comes from
the remaining question in that paper, rather than its already improved
predecessor bounds.

\subsection{Audit --- 2026-09-10}\label{audit-2026-09-10}

Rechecked
\href{https://arxiv.org/html/2608.26047v2}{Lin--Mirrokni--Woodruff,
§1.4} and the August 27 version record. The joint \(k,\varepsilon\) gap
remains explicit. Strong-coreset follow-up searches found no resolution.
Difficulty was raised because the target concerns all input matrices and
requires a new uniform sampling analysis.
\end{document}
